% s-381-p381-m001; original PDF 381
\PoemStart{s-381-p381-m001}{Das Weib.}
\begin{BookVerse}
\VerseLine{0}{Ein Meister sprach zu seinem Gott und Herrn:}
\VerseLine{1}{Nicht einmal nur hast du den Mann beraubt,}
\VerseLine{1}{Da eine Rippe du ihm ausgeklaubt}
\VerseLine{0}{Und Weibsgestalt geschnitzt aus solchem Kern.}
\Stanza
\VerseLine{0}{Sonst strebt der Theil zum Ganzen, fügt sich gern:}
\VerseLine{1}{Doch so verwirrtest Glieder du und Haupt,}
\VerseLine{1}{Dass sich der Mann entwurzelt-hülflos glaubt,}
\VerseLine{0}{Hält ihn vom Weib nur eine Stunde fern.}
\Stanza
\VerseLine{0}{Ich fand kein Weib, das, wie’s das Recht verlangte,}
\VerseLine{0}{Zum Mann zurück, als seinem Ursprung, bangte.}
\VerseLine{1}{Doch fand ich Männer, die nach ihrer Buhle}
\Stanza
\VerseLine{0}{So ernstlich schnappten wie der Fisch nach Wasser.}
\VerseLine{0}{Erfahrung, ach! macht mich zum Weiberhasser:}
\VerseLine{1}{Mein bester Jünger lief mir aus der Schule!}
\end{BookVerse}
\PoemEnd
